\begin{python0}
from solutions import *; clear()
\end{python0}
\sectionthree{What is recursion?}

A \EMPHASIZE{recursive function} is easy to define:
It's a function that calls itself.
(There's another kind of recursion called recursive data
structure.)

Here's an example:

\begin{consolethree}[escapeinside=||]
int f(int x)
{
    std::cout << "entering f() with x = " << x
              << std::endl;
    if (x == 0)
    {
        return 0;
    }
    else
    {
        return f(x - 1);
    }
}

int main()
{
    int x = f(3);
    std::cout << "main() ... x = " << x << std::endl;
    return 0;
}
\end{consolethree}

We say that \texttt{f()} is a \textbf{recursive function}.

Let's trace the program slowly. The difficulty is that
\texttt{f()} is being called several times.

\texttt{main()} calls \texttt{f(3)}

\texttt{f(3)} calls \texttt{f(3 - 1)}, i.e., \texttt{f(2)}

\texttt{f(2)} calls \texttt{f(1)}

\texttt{f(1)} calls \texttt{f(0)}

Note that

\texttt{f(0)} returns 0 to \texttt{f(1)}

\texttt{f(1)} receives 0 and return it to \texttt{f(2)}

\texttt{f(2)} receives 0 and return it to \texttt{f(3)}

\texttt{f(3)} receives 0 and return it to \texttt{main()}

Sometimes it's helpful to visual this stack of function
calls:

\begin{python}
from latextool_basic import *
p = Plot()

p += Rect(0, 0, 2, 1, label=r'\texttt{f(0)}', linewidth=0.05)
p += Rect(0, 1.2, 2, 2.2, label=r'\texttt{f(1)}', linewidth=0.05)
p += Rect(0, 2.4, 2, 3.4, label=r'\texttt{f(2)}', linewidth=0.05)
p += Rect(0, 3.6, 2, 4.6, label=r'\texttt{f(3)}', linewidth=0.05)
p += Rect(0, 4.8, 2, 5.8, label=r'\texttt{main()}', linewidth=0.05)

p += Line(points=[(0, 5.3), (-0.5, 5.3), (-0.5, 4.1), (0, 4.1)], endstyle='>', linewidth=0.05)
p += Line(points=[(0, 3.9), (-0.5, 3.9), (-0.5, 2.9), (0, 2.9)], endstyle='>', linewidth=0.05)
p += Line(points=[(0, 2.7), (-0.5, 2.7), (-0.5, 1.7), (0, 1.7)], endstyle='>', linewidth=0.05)
p += Line(points=[(0, 1.5), (-0.5, 1.5), (-0.5, 0.5), (0, 0.5)], endstyle='>', linewidth=0.05)
# Labels for calls
p += Rect(-1.5, 4.5, -0.5, 5.0, label=r'\small{calls}', linewidth=0)
# Return arrows (right side)
p += Line(points=[(2, 0.5), (2.5, 0.5), (2.5, 1.7), (2, 1.7)], endstyle='>', linewidth=0.05, linestyle='dashed')
p += Line(points=[(2, 1.9), (2.5, 1.9), (2.5, 2.9), (2, 2.9)], endstyle='>', linewidth=0.05, linestyle='dashed')
p += Line(points=[(2, 3.1), (2.5, 3.1), (2.5, 4.1), (2, 4.1)], endstyle='>', linewidth=0.05, linestyle='dashed')
p += Line(points=[(2, 4.3), (2.5, 4.3), (2.5, 5.3), (2, 5.3)], endstyle='>', linewidth=0.05, linestyle='dashed')
# Labels for returns
p += Rect(2.5, 0.7, 4.0, 1.2, label=r'\small{returns 0}', linewidth=0)
p += Rect(2.5, 1.9, 4.0, 2.4, label=r'\small{returns 0}', linewidth=0)
p += Rect(2.5, 3.1, 4.0, 3.6, label=r'\small{returns 0}', linewidth=0)
p += Rect(2.5, 4.3, 4.0, 4.8, label=r'\small{returns 0}', linewidth=0)
print(p)
\end{python}

Our C++ function \texttt{f()}:

\begin{consolethree}[escapeinside=||]
int f(int x)
{
    if (x == 0)
    {
        return 0;
    }
    else
    {
        return f(x - 1);
    }
}
\end{consolethree}

is written mathematically like this:

\[
f(x) = \begin{cases}
0 & \text{if } x = 0 \\
f(x - 1) & \text{if } x > 0
\end{cases}
\]

Now let's do a recursion that does something useful ...


